\documentclass[12pt,a4paper]{article}
\usepackage{xeCJK}
\usepackage{geometry}
\usepackage{graphicx}
\usepackage{hyperref}
\usepackage{listings}
\usepackage{xcolor}
\usepackage{booktabs}
\usepackage{amsmath}

\geometry{left=2.5cm,right=2.5cm,top=2.5cm,bottom=2.5cm}

% 设置中文字体
\setCJKmainfont{STSong}
\setCJKsansfont{STHeiti}
\setCJKmonofont{STFangsong}

% 代码样式
\lstset{
    basicstyle=\ttfamily\small,
    keywordstyle=\color{blue},
    commentstyle=\color{gray},
    stringstyle=\color{red},
    numbers=left,
    numberstyle=\tiny\color{gray},
    frame=single,
    breaklines=true,
    captionpos=b
}

\title{\textbf{CoreNN 文件索引与搜索系统\\技术报告}}
\author{王家齐}
\date{\today}

\begin{document}

\maketitle

\begin{abstract}
本报告详细介绍了基于 CoreNN 向量数据库的大规模文件索引与搜索系统的设计、实现与性能分析。该系统能够对近千万个文件进行高效索引，并提供毫秒级的搜索响应。系统采用 Rust 实现核心向量数据库，使用 Python 构建索引器和 Web 服务，实现了高性能的文件检索功能。
\end{abstract}

\tableofcontents
\newpage

\section{引言}

\subsection{背景}
在现代计算环境中，用户通常需要管理数以百万计的文件。传统的文件搜索方法（如基于文件名的精确匹配）往往无法满足快速、模糊搜索的需求。本项目旨在构建一个基于向量相似度的文件搜索系统，能够快速定位用户所需的文件。

\subsection{目标}
\begin{itemize}
    \item 索引大规模文件系统（支持千万级文件）
    \item 提供毫秒级搜索响应
    \item 支持模糊搜索和语义相似度匹配
    \item 提供友好的 Web 界面和 API 接口
\end{itemize}

\section{系统架构}

\subsection{整体架构}
系统采用三层架构设计：

\begin{enumerate}
    \item \textbf{数据层}：基于 RocksDB 的向量数据库，存储文件向量和元数据
    \item \textbf{服务层}：Python Flask Web 服务，提供搜索 API
    \item \textbf{表示层}：Web 界面，提供用户交互
\end{enumerate}

\subsection{核心组件}

\subsubsection{CoreNN 向量数据库}
CoreNN 是使用 Rust 实现的高性能向量数据库，采用 HNSW (Hierarchical Navigable Small World) 算法进行近似最近邻搜索。

\textbf{主要特性：}
\begin{itemize}
    \item 支持 128 维浮点向量
    \item 基于 RocksDB 的持久化存储
    \item HNSW 索引结构，提供 $O(\log n)$ 搜索复杂度
    \item 支持批量插入和查询
\end{itemize}

\subsubsection{文件索引器}
文件索引器负责扫描文件系统并生成向量表示。

\textbf{向量特征构成（128维）：}
\begin{itemize}
    \item 文件名特征（32维）：基于 MD5 哈希
    \item 扩展名特征（16维）：文件类型标识
    \item 目录路径特征（32维）：文件位置信息
    \item 文件大小特征（16维）：归一化的文件大小
    \item 修改时间特征（16维）：时间戳信息
    \item MIME 类型特征（16维）：文件类型分类
\end{itemize}

\section{实现细节}

\subsection{索引构建流程}

\begin{enumerate}
    \item \textbf{文件扫描}：使用 \texttt{os.walk()} 遍历文件系统
    \item \textbf{特征提取}：为每个文件生成 128 维向量
    \item \textbf{向量归一化}：确保向量模长为 1
    \item \textbf{批量插入}：将向量写入 CoreNN 数据库
    \item \textbf{元数据保存}：将文件信息保存为 JSON
\end{enumerate}

\subsection{搜索算法}

搜索过程分为以下步骤：

\begin{enumerate}
    \item 将查询字符串转换为 128 维向量
    \item 使用 HNSW 算法在向量空间中查找最近邻
    \item 计算余弦相似度：
    \begin{equation}
        \text{similarity} = \frac{\mathbf{q} \cdot \mathbf{v}}{|\mathbf{q}| \cdot |\mathbf{v}|}
    \end{equation}
    \item 返回 Top-K 个最相似的文件
\end{enumerate}

\section{性能分析}

\subsection{索引构建性能}

\begin{table}[h]
\centering
\begin{tabular}{@{}lr@{}}
\toprule
\textbf{指标} & \textbf{数值} \\
\midrule
索引文件总数 & 9,431,950 \\
总耗时 & 2小时10分钟 \\
处理速度 & 72,400 文件/分钟 \\
平均每文件 & 0.83 毫秒 \\
向量数据库大小 & 11 GB \\
元数据大小 & 4.2 GB \\
总索引大小 & 15.2 GB \\
\bottomrule
\end{tabular}
\caption{索引构建性能统计}
\end{table}

\subsection{搜索性能}

\begin{table}[h]
\centering
\begin{tabular}{@{}lc@{}}
\toprule
\textbf{操作} & \textbf{响应时间} \\
\midrule
单次查询 & < 100 ms \\
Top-10 结果 & 50-80 ms \\
Top-50 结果 & 80-150 ms \\
\bottomrule
\end{tabular}
\caption{搜索响应时间}
\end{table}

\subsection{内存使用}

\begin{itemize}
    \item 索引构建峰值内存：25.8 GB
    \item 搜索服务运行内存：20.3 GB
    \item 向量数据库加载时间：约 30 秒
\end{itemize}

\section{系统使用}

\subsection{启动服务}

\begin{lstlisting}[language=bash,caption=启动搜索服务器]
# 激活虚拟环境
source venv/bin/activate

# 启动服务器
python file_search_server.py
\end{lstlisting}

\subsection{API 使用示例}

\begin{lstlisting}[language=bash,caption=搜索 API 调用]
# 搜索文件
curl -X POST http://localhost:5000/search \
  -H "Content-Type: application/json" \
  -d '{"query": "config.json", "limit": 10}'

# 查看统计信息
curl http://localhost:5000/stats
\end{lstlisting}

\section{问题与优化}

\subsection{当前问题}

\subsubsection{搜索结果准确性}
当前搜索算法基于简单的 MD5 哈希特征，可能导致以下问题：
\begin{itemize}
    \item 语义相似但名称不同的文件无法关联
    \item 哈希冲突可能影响搜索准确性
    \item 缺乏对文件内容的理解
\end{itemize}

\subsection{优化方案}

\subsubsection{改进向量表示}
\begin{enumerate}
    \item \textbf{使用 N-gram 特征}：将文件名分解为字符级 N-gram，提高模糊匹配能力
    \item \textbf{TF-IDF 权重}：对常见词降权，提高区分度
    \item \textbf{语义嵌入}：使用预训练模型（如 BERT）生成语义向量
\end{enumerate}

\subsubsection{优化搜索算法}
\begin{enumerate}
    \item \textbf{多阶段检索}：先粗筛选，再精排序
    \item \textbf{查询扩展}：自动添加同义词和相关词
    \item \textbf{个性化排序}：根据用户历史调整结果
\end{enumerate}

\section{结论}

本项目成功实现了一个高性能的大规模文件索引与搜索系统，能够在毫秒级时间内从近千万个文件中检索目标文件。系统采用向量相似度搜索技术，提供了比传统文件搜索更灵活的检索方式。

\subsection{主要成果}
\begin{itemize}
    \item 成功索引 9,431,950 个文件
    \item 实现毫秒级搜索响应（< 100ms）
    \item 提供 Web 界面和 API 接口
    \item 系统稳定运行，内存占用合理
\end{itemize}

\subsection{未来工作}
\begin{itemize}
    \item 改进向量表示方法，提高搜索准确性
    \item 支持增量索引更新
    \item 添加文件内容索引功能
    \item 优化内存使用，降低系统资源消耗
\end{itemize}

\section*{参考文献}
\begin{enumerate}
    \item Malkov, Y. A., \& Yashunin, D. A. (2018). Efficient and robust approximate nearest neighbor search using hierarchical navigable small world graphs. IEEE transactions on pattern analysis and machine intelligence, 42(4), 824-836.
    \item RocksDB: A Persistent Key-Value Store for Fast Storage Environments. Facebook, Inc.
    \item Johnson, J., Douze, M., \& Jégou, H. (2019). Billion-scale similarity search with GPUs. IEEE Transactions on Big Data.
\end{enumerate}

\end{document}
